Agentic LLM workloads expand a committed context into $k$ candidate continuations in a Tree-of-Thoughts style fan-out, of which at most one spine is retained.
Under prefill--decode (PD) disaggregation, each candidate is prefilled and its KV cache transferred before decode-side path pruning can take effect, so non-surviving branches still pay residual prefill and connector cost that later Tree-of-Thoughts pruning cannot recover.
We propose \emph{PD pruning disaggregation}, which coordinates two complementary forms of pruning across the disaggregated path: prefill-side admission over the shared prefix tree, and decode-side exploration over the reasoning tree.
On the prefill side, \sys selectively admits residuals into the radix-style KV tree before kernels run---reusing resident trunk pages, gating which suffixes are computed, and transferring only pages that the decode instance still lacks---so unadmitted candidates never consume residual prefill or connector bandwidth.
On the decode side, candidates whose continuation remains unresolved receive a short probe and are expanded further only when that probe supports continuation, while Tree-of-Thoughts path pruners such as ESC, Speculative Rejection, and DPTS operate unchanged on the admitted set; a copy-on-write context tree then reclaims aborted residual pages and retains the committed spine as the next trunk.
On Qwen3-8B, prefill-side admission reduces computed prefill by $23\%$ while maintaining reasoning quality; on MATH-500, decode-side probing accelerates serving to $0.70\times$ the wall-clock of the full fan-out while maintaining coverage ($20/32$); on GSM8K with Qwen3-14B, the retained spine requires $0.72\times$ the peak KV of prefix caching while maintaining reasoning quality.
